%%%%%%%%%%%%%%%%%%%%%%%%%%%%%%%%%%%%%%%%%%%%%%%%%%%%%%%%%%%%
% 6.7 ALGEBRAIC COMBINATORICS
% The Composition of Advanced Mathematics
%%%%%%%%%%%%%%%%%%%%%%%%%%%%%%%%%%%%%%%%%%%%%%%%%%%%%%%%%%%%

\section{Algebraic Combinatorics}
\label{sec:algebraic-combinatorics}

\prerequisite{
Permutations, group actions, basic linear algebra, polynomial algebra,
generating functions, integer partitions, and the enumerative methods
developed earlier in this chapter.
}

\learningobjectives{
By the end of this section, the reader should be able to:
\begin{itemize}
    \item explain the role of algebra in combinatorial enumeration;
    \item use group actions to formalize symmetry;
    \item compute orbits and stabilizers;
    \item apply the orbit--stabilizer theorem;
    \item apply Burnside's lemma to count objects up to symmetry;
    \item use cycle structure in permutation actions;
    \item understand the purpose of P\'olya enumeration;
    \item work with cycle index polynomials;
    \item distinguish labeled and unlabeled enumeration;
    \item understand permutation statistics algebraically;
    \item define \(q\)-integers, \(q\)-factorials, and Gaussian binomial
          coefficients;
    \item interpret Gaussian coefficients as finite-field subspace
          counts;
    \item understand Young diagrams and standard Young tableaux;
    \item apply the hook-length formula in basic cases;
    \item recognize symmetric functions as generating objects;
    \item understand the roles of Schur functions and representation
          theory;
    \item recognize association schemes and algebraic graph invariants;
    \item connect algebraic structure with combinatorial classification
          and enumeration.
\end{itemize}
}

%%%%%%%%%%%%%%%%%%%%%%%%%%%%%%%%%%%%%%%%%%%%%%%%%%%%%%%%%%%%
\subsection{What Is Algebraic Combinatorics?}
\label{subsec:what-is-algebraic-combinatorics}
%%%%%%%%%%%%%%%%%%%%%%%%%%%%%%%%%%%%%%%%%%%%%%%%%%%%%%%%%%%%

Algebraic combinatorics uses algebraic structures and techniques to
study discrete objects.

The central philosophy is

\[
\boxed{
\text{combinatorial structure}
\longrightarrow
\text{algebraic representation}
\longrightarrow
\text{structural or enumerative information}.
}
\]

Typical algebraic tools include:

\begin{itemize}
    \item groups and group actions;
    \item rings and polynomial algebras;
    \item vector spaces;
    \item matrices;
    \item symmetric functions;
    \item representation theory;
    \item finite fields.
\end{itemize}

These tools frequently reveal symmetries and identities that are
difficult to see by direct counting alone.

%%%%%%%%%%%%%%%%%%%%%%%%%%%%%%%%%%%%%%%%%%%%%%%%%%%%%%%%%%%%
\subsection{Symmetry and Group Actions}
\label{subsec:symmetry-group-actions}
%%%%%%%%%%%%%%%%%%%%%%%%%%%%%%%%%%%%%%%%%%%%%%%%%%%%%%%%%%%%

Suppose a group \(G\) acts on a set \(X\).

For

\[
g\in G
\]

and

\[
x\in X,
\]

the image of \(x\) under \(g\) is written

\[
g\cdot x.
\]

A group action satisfies

\[
\boxed{
e\cdot x=x
}
\]

for the identity element \(e\), and

\[
\boxed{
g\cdot(h\cdot x)
=
(gh)\cdot x.
}
\]

Group actions provide the natural framework for describing symmetry.

%%%%%%%%%%%%%%%%%%%%%%%%%%%%%%%%%%%%%%%%%%%%%%%%%%%%%%%%%%%%
\subsection{Orbits}
\label{subsec:orbits-algebraic-combinatorics}
%%%%%%%%%%%%%%%%%%%%%%%%%%%%%%%%%%%%%%%%%%%%%%%%%%%%%%%%%%%%

\begin{definitionbox}{Orbit}
The orbit of \(x\in X\) under \(G\) is

\[
\boxed{
\operatorname{Orb}(x)
=
\{g\cdot x:g\in G\}.
}
\]
\end{definitionbox}

Two objects lie in the same orbit exactly when one can be transformed
into the other by an allowed symmetry.

Thus orbits represent equivalence classes under symmetry.

%%%%%%%%%%%%%%%%%%%%%%%%%%%%%%%%%%%%%%%%%%%%%%%%%%%%%%%%%%%%
\subsection{Stabilizers}
\label{subsec:stabilizers-algebraic-combinatorics}
%%%%%%%%%%%%%%%%%%%%%%%%%%%%%%%%%%%%%%%%%%%%%%%%%%%%%%%%%%%%

\begin{definitionbox}{Stabilizer}
The stabilizer of \(x\in X\) is

\[
\boxed{
\operatorname{Stab}(x)
=
\{g\in G:g\cdot x=x\}.
}
\]
\end{definitionbox}

The stabilizer consists of the symmetries that leave \(x\) unchanged.

Highly symmetric objects have large stabilizers.

Generic objects often have smaller stabilizers.

%%%%%%%%%%%%%%%%%%%%%%%%%%%%%%%%%%%%%%%%%%%%%%%%%%%%%%%%%%%%
\subsection{Orbit--Stabilizer Theorem}
\label{subsec:orbit-stabilizer}
%%%%%%%%%%%%%%%%%%%%%%%%%%%%%%%%%%%%%%%%%%%%%%%%%%%%%%%%%%%%

\begin{theorem}[Orbit--Stabilizer Theorem]
If a finite group \(G\) acts on \(X\), then for every \(x\in X\),

\[
\boxed{
|G|
=
|\operatorname{Orb}(x)|
|\operatorname{Stab}(x)|.
}
\]
\end{theorem}

Equivalently,

\[
\boxed{
|\operatorname{Orb}(x)|
=
\frac{|G|}{|\operatorname{Stab}(x)|}.
}
\]

This explains why symmetric objects may have smaller orbits than
generic objects.

%%%%%%%%%%%%%%%%%%%%%%%%%%%%%%%%%%%%%%%%%%%%%%%%%%%%%%%%%%%%
\subsection{Worked Example: Rotating a Square}
\label{subsec:worked-square-orbit}
%%%%%%%%%%%%%%%%%%%%%%%%%%%%%%%%%%%%%%%%%%%%%%%%%%%%%%%%%%%%

\begin{workedexamplebox}{Orbit of a Marked Vertex}

Consider the rotational symmetries of a square.

The rotation group has size

\[
|G|=4.
\]

Mark one vertex.

Only the identity rotation fixes that marked vertex, so

\[
|\operatorname{Stab}(x)|=1.
\]

Therefore,

\[
|\operatorname{Orb}(x)|
=
\frac41
=
4.
\]

\begin{answerbox}
\[
\boxed{
|\operatorname{Orb}(x)|=4.
}
\]
\end{answerbox}

\end{workedexamplebox}

%%%%%%%%%%%%%%%%%%%%%%%%%%%%%%%%%%%%%%%%%%%%%%%%%%%%%%%%%%%%
\subsection{Counting Up to Symmetry}
\label{subsec:counting-up-to-symmetry}
%%%%%%%%%%%%%%%%%%%%%%%%%%%%%%%%%%%%%%%%%%%%%%%%%%%%%%%%%%%%

Suppose one wants to count colorings of a geometric object.

If rotations or reflections are considered equivalent, ordinary
counting overcounts.

The correct objects are the orbits of the symmetry group.

Thus the problem becomes

\[
\boxed{
\text{count the orbits of }G\text{ acting on }X.
}
\]

Burnside's lemma gives a direct formula.

%%%%%%%%%%%%%%%%%%%%%%%%%%%%%%%%%%%%%%%%%%%%%%%%%%%%%%%%%%%%
\subsection{Fixed Points of a Group Element}
\label{subsec:fixed-points-group-action}
%%%%%%%%%%%%%%%%%%%%%%%%%%%%%%%%%%%%%%%%%%%%%%%%%%%%%%%%%%%%

For

\[
g\in G,
\]

define

\[
\boxed{
\operatorname{Fix}(g)
=
\{x\in X:g\cdot x=x\}.
}
\]

The quantity

\[
|\operatorname{Fix}(g)|
\]

counts the objects unchanged by the symmetry \(g\).

Different symmetries generally fix different numbers of objects.

%%%%%%%%%%%%%%%%%%%%%%%%%%%%%%%%%%%%%%%%%%%%%%%%%%%%%%%%%%%%
\subsection{Burnside's Lemma}
\label{subsec:burnsides-lemma}
%%%%%%%%%%%%%%%%%%%%%%%%%%%%%%%%%%%%%%%%%%%%%%%%%%%%%%%%%%%%

\begin{theorem}[Burnside's Lemma]
If a finite group \(G\) acts on a finite set \(X\), then the number of
orbits is

\[
\boxed{
|X/G|
=
\frac{1}{|G|}
\sum_{g\in G}
|\operatorname{Fix}(g)|.
}
\]
\end{theorem}

This is also called the Cauchy--Frobenius lemma.

The notation

\[
X/G
\]

denotes the set of orbits.

%%%%%%%%%%%%%%%%%%%%%%%%%%%%%%%%%%%%%%%%%%%%%%%%%%%%%%%%%%%%
\subsection{Why Burnside's Lemma Works}
\label{subsec:why-burnside-works}
%%%%%%%%%%%%%%%%%%%%%%%%%%%%%%%%%%%%%%%%%%%%%%%%%%%%%%%%%%%%

Count pairs

\[
(g,x)
\]

such that

\[
g\cdot x=x.
\]

Counting by \(g\) gives

\[
\sum_{g\in G}
|\operatorname{Fix}(g)|.
\]

Counting by \(x\) gives

\[
\sum_{x\in X}
|\operatorname{Stab}(x)|.
\]

Group the second sum by orbits.

For each orbit \(\mathcal{O}\),

\[
|\mathcal{O}|
|\operatorname{Stab}(x)|
=
|G|.
\]

Thus every orbit contributes exactly \(|G|\).

Hence,

\[
\sum_{g\in G}
|\operatorname{Fix}(g)|
=
|G|\cdot|X/G|.
\]

Dividing by \(|G|\) gives Burnside's formula.

%%%%%%%%%%%%%%%%%%%%%%%%%%%%%%%%%%%%%%%%%%%%%%%%%%%%%%%%%%%%
\subsection{Worked Example: Two-Color Necklaces}
\label{subsec:worked-burnside-necklace}
%%%%%%%%%%%%%%%%%%%%%%%%%%%%%%%%%%%%%%%%%%%%%%%%%%%%%%%%%%%%

\begin{workedexamplebox}{Coloring Four Beads Up to Rotation}

Consider four positions arranged in a circle.

Each position may be colored black or white.

There are

\[
2^4=16
\]

labeled colorings.

Let the cyclic rotation group \(C_4\) act.

The identity fixes all

\[
16
\]

colorings.

A \(90^\circ\) rotation fixes only colorings in which all four positions
have the same color:

\[
2.
\]

A \(180^\circ\) rotation has two cycles of positions and therefore fixes

\[
2^2=4
\]

colorings.

A \(270^\circ\) rotation again fixes

\[
2
\]

colorings.

Burnside's lemma gives

\[
\frac{16+2+4+2}{4}
=
6.
\]

\begin{answerbox}
\[
\boxed{
6
}
\]

binary necklaces of length \(4\) exist up to rotation.
\end{answerbox}

\end{workedexamplebox}

%%%%%%%%%%%%%%%%%%%%%%%%%%%%%%%%%%%%%%%%%%%%%%%%%%%%%%%%%%%%
\subsection{Cycle Structure and Fixed Colorings}
\label{subsec:cycle-structure-fixed-colorings}
%%%%%%%%%%%%%%%%%%%%%%%%%%%%%%%%%%%%%%%%%%%%%%%%%%%%%%%%%%%%

Suppose a permutation \(g\) acts on \(n\) positions.

If the action decomposes the positions into

\[
c(g)
\]

cycles, then a coloring fixed by \(g\) must be constant on every cycle.

With \(q\) available colors,

\[
\boxed{
|\operatorname{Fix}(g)|
=
q^{c(g)}.
}
\]

This simple observation is the basis of many symmetry-counting formulas.

%%%%%%%%%%%%%%%%%%%%%%%%%%%%%%%%%%%%%%%%%%%%%%%%%%%%%%%%%%%%
\subsection{Necklaces}
\label{subsec:necklaces}
%%%%%%%%%%%%%%%%%%%%%%%%%%%%%%%%%%%%%%%%%%%%%%%%%%%%%%%%%%%%

A necklace is a circular arrangement of colored beads in which
rotations are considered equivalent.

For length \(n\) with \(q\) colors, Burnside's lemma gives

\[
\boxed{
N_q(n)
=
\frac1n
\sum_{k=0}^{n-1}
q^{\gcd(n,k)}.
}
\]

The notation

\[
\gcd(n,k)
\]

means the greatest common divisor of \(n\) and \(k\).

An equivalent divisor-sum formula is

\[
\boxed{
N_q(n)
=
\frac1n
\sum_{d\mid n}
\varphi(d)\,
q^{n/d},
}
\]

where \(\varphi\) is Euler's totient function.

%%%%%%%%%%%%%%%%%%%%%%%%%%%%%%%%%%%%%%%%%%%%%%%%%%%%%%%%%%%%
\subsection{Bracelets}
\label{subsec:bracelets}
%%%%%%%%%%%%%%%%%%%%%%%%%%%%%%%%%%%%%%%%%%%%%%%%%%%%%%%%%%%%

A bracelet identifies colorings under both rotations and reflections.

The relevant symmetry group is the dihedral group

\[
D_n.
\]

Thus bracelets are counted using Burnside's lemma over all:

\begin{itemize}
    \item rotations;
    \item reflections.
\end{itemize}

The reflection contribution depends on whether \(n\) is even or odd.

This illustrates how the group itself determines the counting formula.

%%%%%%%%%%%%%%%%%%%%%%%%%%%%%%%%%%%%%%%%%%%%%%%%%%%%%%%%%%%%
\subsection{Cycle Index}
\label{subsec:cycle-index}
%%%%%%%%%%%%%%%%%%%%%%%%%%%%%%%%%%%%%%%%%%%%%%%%%%%%%%%%%%%%

Burnside's lemma can be organized using a cycle index polynomial.

Suppose \(G\) acts on \(n\) positions.

For \(g\in G\), let

\[
c_i(g)
\]

be the number of cycles of length \(i\) in the permutation induced by
\(g\).

\begin{definitionbox}{Cycle Index Polynomial}
The cycle index of \(G\) is

\[
\boxed{
Z_G(a_1,\ldots,a_n)
=
\frac{1}{|G|}
\sum_{g\in G}
\prod_{i=1}^{n}
a_i^{c_i(g)}.
}
\]
\end{definitionbox}

The variables

\[
a_1,a_2,\ldots
\]

record cycle lengths.

%%%%%%%%%%%%%%%%%%%%%%%%%%%%%%%%%%%%%%%%%%%%%%%%%%%%%%%%%%%%
\subsection{Cycle Index of a Cyclic Group}
\label{subsec:cycle-index-cyclic}
%%%%%%%%%%%%%%%%%%%%%%%%%%%%%%%%%%%%%%%%%%%%%%%%%%%%%%%%%%%%

For the cyclic group acting on \(n\) positions,

\[
C_n,
\]

the cycle index can be written

\[
\boxed{
Z_{C_n}
=
\frac1n
\sum_{d\mid n}
\varphi(d)
a_d^{\,n/d}.
}
\]

This compact expression organizes the rotational symmetry of an
\(n\)-cycle.

%%%%%%%%%%%%%%%%%%%%%%%%%%%%%%%%%%%%%%%%%%%%%%%%%%%%%%%%%%%%
\subsection{P\'olya Enumeration}
\label{subsec:polya-enumeration}
%%%%%%%%%%%%%%%%%%%%%%%%%%%%%%%%%%%%%%%%%%%%%%%%%%%%%%%%%%%%

P\'olya enumeration generalizes Burnside's lemma by allowing different
colors to carry weights.

Suppose colors have weights

\[
w_1,w_2,\ldots,w_q.
\]

In the cycle index, substitute

\[
\boxed{
a_i
\longmapsto
w_1^i+w_2^i+\cdots+w_q^i.
}
\]

The resulting polynomial records colorings up to symmetry according to
their color content.

%%%%%%%%%%%%%%%%%%%%%%%%%%%%%%%%%%%%%%%%%%%%%%%%%%%%%%%%%%%%
\subsection{Worked Example: Cycle Index of \(C_4\)}
\label{subsec:worked-cycle-index-c4}
%%%%%%%%%%%%%%%%%%%%%%%%%%%%%%%%%%%%%%%%%%%%%%%%%%%%%%%%%%%%

\begin{workedexamplebox}{Cycle Index for Four Rotational Positions}

For the rotations of four positions:

\begin{itemize}
    \item the identity has cycle type \(1^4\);
    \item the \(90^\circ\) and \(270^\circ\) rotations each have cycle
          type \(4^1\);
    \item the \(180^\circ\) rotation has cycle type \(2^2\).
\end{itemize}

Therefore,

\[
Z_{C_4}
=
\frac14
\left(
a_1^4
+
a_2^2
+
2a_4
\right).
\]

\begin{answerbox}
\[
\boxed{
Z_{C_4}
=
\frac14
\left(
a_1^4+a_2^2+2a_4
\right).
}
\]
\end{answerbox}

\end{workedexamplebox}

%%%%%%%%%%%%%%%%%%%%%%%%%%%%%%%%%%%%%%%%%%%%%%%%%%%%%%%%%%%%
\subsection{P\'olya Substitution Example}
\label{subsec:polya-substitution-example}
%%%%%%%%%%%%%%%%%%%%%%%%%%%%%%%%%%%%%%%%%%%%%%%%%%%%%%%%%%%%

For two colors, say \(x\) and \(y\), substitute

\[
a_i
\longmapsto
x^i+y^i.
\]

For \(C_4\),

\[
Z_{C_4}
=
\frac14
\left(
a_1^4+a_2^2+2a_4
\right)
\]

becomes

\[
\boxed{
\frac14
\left[
(x+y)^4
+
(x^2+y^2)^2
+
2(x^4+y^4)
\right].
}
\]

The coefficient of

\[
x^ry^{4-r}
\]

counts rotational equivalence classes using exactly \(r\) objects of the
first color.

%%%%%%%%%%%%%%%%%%%%%%%%%%%%%%%%%%%%%%%%%%%%%%%%%%%%%%%%%%%%
\subsection{Symmetry Counting and Chemical Isomers}
\label{subsec:symmetry-chemical-isomers}
%%%%%%%%%%%%%%%%%%%%%%%%%%%%%%%%%%%%%%%%%%%%%%%%%%%%%%%%%%%%

P\'olya enumeration was historically important in counting chemical
isomers.

If positions on a molecular structure are permuted by geometric
symmetries, different substitutions may represent the same molecule.

Thus one studies

\[
\boxed{
\text{colorings of positions}
\big/
\text{molecular symmetry group}.
}
\]

The same mathematics appears in:

\begin{itemize}
    \item necklaces;
    \item graph colorings;
    \item molecular substitution patterns;
    \item tilings;
    \item crystal structures.
\end{itemize}

%%%%%%%%%%%%%%%%%%%%%%%%%%%%%%%%%%%%%%%%%%%%%%%%%%%%%%%%%%%%
\subsection{Permutation Statistics}
\label{subsec:algebraic-permutation-statistics}
%%%%%%%%%%%%%%%%%%%%%%%%%%%%%%%%%%%%%%%%%%%%%%%%%%%%%%%%%%%%

Algebraic combinatorics often refines permutation enumeration using
statistics.

For

\[
\pi\in S_n,
\]

examples include

\[
\operatorname{inv}(\pi)
\]

for inversion number,

\[
\operatorname{des}(\pi)
\]

for number of descents,

and

\[
\operatorname{cyc}(\pi)
\]

for number of cycles.

One may form generating polynomials such as

\[
\boxed{
\sum_{\pi\in S_n}
q^{\operatorname{inv}(\pi)}.
}
\]

%%%%%%%%%%%%%%%%%%%%%%%%%%%%%%%%%%%%%%%%%%%%%%%%%%%%%%%%%%%%
\subsection{\(q\)-Integers}
\label{subsec:q-integers-algebraic}
%%%%%%%%%%%%%%%%%%%%%%%%%%%%%%%%%%%%%%%%%%%%%%%%%%%%%%%%%%%%

Define

\[
\boxed{
[n]_q
=
1+q+q^2+\cdots+q^{n-1}.
}
\]

When \(q\neq1\),

\[
\boxed{
[n]_q
=
\frac{1-q^n}{1-q}.
}
\]

As \(q\to1\),

\[
[n]_q\to n.
\]

Thus \([n]_q\) is called a \(q\)-analogue of \(n\).

%%%%%%%%%%%%%%%%%%%%%%%%%%%%%%%%%%%%%%%%%%%%%%%%%%%%%%%%%%%%
\subsection{\(q\)-Factorials}
\label{subsec:q-factorials-algebraic}
%%%%%%%%%%%%%%%%%%%%%%%%%%%%%%%%%%%%%%%%%%%%%%%%%%%%%%%%%%%%

Define

\[
\boxed{
[n]_q!
=
[1]_q[2]_q\cdots[n]_q.
}
\]

One important identity is

\[
\boxed{
[n]_q!
=
\sum_{\pi\in S_n}
q^{\operatorname{inv}(\pi)}.
}
\]

Thus the \(q\)-factorial refines \(n!\) by tracking inversion number.

Setting

\[
q=1
\]

recovers

\[
n!.
\]

%%%%%%%%%%%%%%%%%%%%%%%%%%%%%%%%%%%%%%%%%%%%%%%%%%%%%%%%%%%%
\subsection{Gaussian Binomial Coefficients}
\label{subsec:gaussian-binomial-coefficients}
%%%%%%%%%%%%%%%%%%%%%%%%%%%%%%%%%%%%%%%%%%%%%%%%%%%%%%%%%%%%

\begin{definitionbox}{Gaussian Binomial Coefficient}
The Gaussian binomial coefficient is

\[
\boxed{
\begin{bmatrix}
n\\
k
\end{bmatrix}_q
=
\frac{
[n]_q!
}{
[k]_q![n-k]_q!
}.
}
\]
\end{definitionbox}

It is also called the \(q\)-binomial coefficient.

As \(q\to1\),

\[
\boxed{
\begin{bmatrix}
n\\
k
\end{bmatrix}_q
\longrightarrow
\binom{n}{k}.
}
\]

%%%%%%%%%%%%%%%%%%%%%%%%%%%%%%%%%%%%%%%%%%%%%%%%%%%%%%%%%%%%
\subsection{Finite-Field Interpretation}
\label{subsec:gaussian-finite-field}
%%%%%%%%%%%%%%%%%%%%%%%%%%%%%%%%%%%%%%%%%%%%%%%%%%%%%%%%%%%%

Let \(q\) be a prime power.

Then the number of \(k\)-dimensional subspaces of the vector space

\[
\mathbb{F}_q^n
\]

is

\[
\boxed{
\begin{bmatrix}
n\\
k
\end{bmatrix}_q.
}
\]

Thus ordinary binomial coefficients count subsets, while Gaussian
binomial coefficients count subspaces.

The analogy is

\[
\boxed{
\text{subsets of an }n\text{-set}
\quad\leftrightarrow\quad
\text{subspaces of }\mathbb{F}_q^n.
}
\]

%%%%%%%%%%%%%%%%%%%%%%%%%%%%%%%%%%%%%%%%%%%%%%%%%%%%%%%%%%%%
\subsection{Worked Example: Lines in \(\mathbb{F}_q^2\)}
\label{subsec:worked-gaussian-lines}
%%%%%%%%%%%%%%%%%%%%%%%%%%%%%%%%%%%%%%%%%%%%%%%%%%%%%%%%%%%%

\begin{workedexamplebox}{One-Dimensional Subspaces}

The number of one-dimensional subspaces of

\[
\mathbb{F}_q^2
\]

is

\[
\begin{bmatrix}
2\\
1
\end{bmatrix}_q.
\]

Using

\[
[2]_q=1+q,
\]

we obtain

\[
\begin{bmatrix}
2\\
1
\end{bmatrix}_q
=
1+q.
\]

\begin{answerbox}
\[
\boxed{
q+1
}
\]

one-dimensional subspaces exist.
\end{answerbox}

\end{workedexamplebox}

%%%%%%%%%%%%%%%%%%%%%%%%%%%%%%%%%%%%%%%%%%%%%%%%%%%%%%%%%%%%
\subsection{\(q\)-Binomial Recurrence}
\label{subsec:q-binomial-recurrence}
%%%%%%%%%%%%%%%%%%%%%%%%%%%%%%%%%%%%%%%%%%%%%%%%%%%%%%%%%%%%

Gaussian coefficients satisfy a \(q\)-analogue of Pascal's identity:

\[
\boxed{
\begin{bmatrix}
n\\
k
\end{bmatrix}_q
=
\begin{bmatrix}
n-1\\
k
\end{bmatrix}_q
+
q^{n-k}
\begin{bmatrix}
n-1\\
k-1
\end{bmatrix}_q.
}
\]

Another equivalent version is

\[
\boxed{
\begin{bmatrix}
n\\
k
\end{bmatrix}_q
=
q^k
\begin{bmatrix}
n-1\\
k
\end{bmatrix}_q
+
\begin{bmatrix}
n-1\\
k-1
\end{bmatrix}_q.
}
\]

%%%%%%%%%%%%%%%%%%%%%%%%%%%%%%%%%%%%%%%%%%%%%%%%%%%%%%%%%%%%
\subsection{Young Diagrams}
\label{subsec:young-diagrams-algebraic}
%%%%%%%%%%%%%%%%%%%%%%%%%%%%%%%%%%%%%%%%%%%%%%%%%%%%%%%%%%%%

Let

\[
\lambda
=
(\lambda_1,\lambda_2,\ldots,\lambda_r)
\]

be a partition of \(n\), written

\[
\lambda\vdash n.
\]

A Young diagram represents the partition using left-aligned rows of
boxes whose lengths are

\[
\lambda_1,\lambda_2,\ldots,\lambda_r.
\]

For example,

\[
\lambda=(4,2,1)
\]

has seven boxes.

Young diagrams encode integer partitions geometrically.

%%%%%%%%%%%%%%%%%%%%%%%%%%%%%%%%%%%%%%%%%%%%%%%%%%%%%%%%%%%%
\subsection{Young Tableaux}
\label{subsec:young-tableaux-algebraic}
%%%%%%%%%%%%%%%%%%%%%%%%%%%%%%%%%%%%%%%%%%%%%%%%%%%%%%%%%%%%

A Young tableau is obtained by filling a Young diagram with entries.

A standard Young tableau of size \(n\):

\begin{itemize}
    \item uses each of
    \[
    1,2,\ldots,n
    \]
    exactly once;
    \item increases across each row;
    \item increases down each column.
\end{itemize}

Standard Young tableaux play a central role in the representation theory
of symmetric groups.

%%%%%%%%%%%%%%%%%%%%%%%%%%%%%%%%%%%%%%%%%%%%%%%%%%%%%%%%%%%%
\subsection{Hook Length}
\label{subsec:hook-length}
%%%%%%%%%%%%%%%%%%%%%%%%%%%%%%%%%%%%%%%%%%%%%%%%%%%%%%%%%%%%

For a cell \(u\) in a Young diagram, its hook consists of:

\begin{itemize}
    \item the cell \(u\);
    \item all cells directly to its right;
    \item all cells directly below it.
\end{itemize}

The number of cells in this hook is the hook length

\[
\boxed{
h(u).
}
\]

%%%%%%%%%%%%%%%%%%%%%%%%%%%%%%%%%%%%%%%%%%%%%%%%%%%%%%%%%%%%
\subsection{Hook-Length Formula}
\label{subsec:hook-length-formula}
%%%%%%%%%%%%%%%%%%%%%%%%%%%%%%%%%%%%%%%%%%%%%%%%%%%%%%%%%%%%

\begin{theorem}[Hook-Length Formula]
If

\[
\lambda\vdash n,
\]

then the number \(f^\lambda\) of standard Young tableaux of shape
\(\lambda\) is

\[
\boxed{
f^\lambda
=
\frac{n!}{
\prod_{u\in\lambda}h(u)
}.
}
\]
\end{theorem}

The product runs over every cell in the Young diagram.

%%%%%%%%%%%%%%%%%%%%%%%%%%%%%%%%%%%%%%%%%%%%%%%%%%%%%%%%%%%%
\subsection{Worked Example: Hook-Length Formula}
\label{subsec:worked-hook-length}
%%%%%%%%%%%%%%%%%%%%%%%%%%%%%%%%%%%%%%%%%%%%%%%%%%%%%%%%%%%%

\begin{workedexamplebox}{Shape \((2,1)\)}

Consider the partition

\[
\lambda=(2,1).
\]

There are \(3\) cells.

Their hook lengths are

\[
3,\qquad1,\qquad1.
\]

Therefore,

\[
f^{(2,1)}
=
\frac{3!}{3\cdot1\cdot1}.
\]

\begin{answerbox}
\[
\boxed{
f^{(2,1)}=2.
}
\]
\end{answerbox}

\end{workedexamplebox}

%%%%%%%%%%%%%%%%%%%%%%%%%%%%%%%%%%%%%%%%%%%%%%%%%%%%%%%%%%%%
\subsection{Young's Lattice}
\label{subsec:youngs-lattice}
%%%%%%%%%%%%%%%%%%%%%%%%%%%%%%%%%%%%%%%%%%%%%%%%%%%%%%%%%%%%

Integer partitions may themselves be organized into a partially ordered
set.

Define

\[
\lambda\leq\mu
\]

when the Young diagram of \(\lambda\) is contained inside the Young
diagram of \(\mu\).

The resulting poset is called Young's lattice.

A covering relation corresponds to adding a single box.

Thus algebraic combinatorics links:

\[
\boxed{
\text{partitions}
+
\text{posets}
+
\text{tableaux}.
}
\]

%%%%%%%%%%%%%%%%%%%%%%%%%%%%%%%%%%%%%%%%%%%%%%%%%%%%%%%%%%%%
\subsection{Symmetric Polynomials}
\label{subsec:symmetric-polynomials}
%%%%%%%%%%%%%%%%%%%%%%%%%%%%%%%%%%%%%%%%%%%%%%%%%%%%%%%%%%%%

A polynomial

\[
f(x_1,\ldots,x_n)
\]

is symmetric if permuting the variables leaves the polynomial unchanged.

For example,

\[
x_1+x_2+\cdots+x_n
\]

is symmetric.

So is

\[
\sum_{i<j}x_ix_j.
\]

Symmetric polynomials naturally encode objects whose labels may be
permuted.

%%%%%%%%%%%%%%%%%%%%%%%%%%%%%%%%%%%%%%%%%%%%%%%%%%%%%%%%%%%%
\subsection{Elementary Symmetric Functions}
\label{subsec:elementary-symmetric-functions}
%%%%%%%%%%%%%%%%%%%%%%%%%%%%%%%%%%%%%%%%%%%%%%%%%%%%%%%%%%%%

The \(k\)-th elementary symmetric polynomial is

\[
\boxed{
e_k(x_1,\ldots,x_n)
=
\sum_{1\leq i_1<\cdots<i_k\leq n}
x_{i_1}\cdots x_{i_k}.
}
\]

For example,

\[
e_1=x_1+\cdots+x_n,
\]

and

\[
e_2
=
\sum_{i<j}x_ix_j.
\]

These functions appear as coefficients of

\[
\boxed{
\prod_{i=1}^{n}(1+x_it)
=
\sum_{k=0}^{n}
e_k t^k.
}
\]

%%%%%%%%%%%%%%%%%%%%%%%%%%%%%%%%%%%%%%%%%%%%%%%%%%%%%%%%%%%%
\subsection{Complete Homogeneous Symmetric Functions}
\label{subsec:complete-homogeneous-symmetric}
%%%%%%%%%%%%%%%%%%%%%%%%%%%%%%%%%%%%%%%%%%%%%%%%%%%%%%%%%%%%

The complete homogeneous symmetric function \(h_k\) is the sum of all
monomials of total degree \(k\):

\[
\boxed{
h_k(x_1,\ldots,x_n)
=
\sum_{i_1+\cdots+i_n=k}
x_1^{i_1}\cdots x_n^{i_n}.
}
\]

Its generating function is

\[
\boxed{
\sum_{k=0}^{\infty}
h_k t^k
=
\prod_{i=1}^{n}
\frac{1}{1-x_it}.
}
\]

Thus symmetric functions connect directly with generating functions.

%%%%%%%%%%%%%%%%%%%%%%%%%%%%%%%%%%%%%%%%%%%%%%%%%%%%%%%%%%%%
\subsection{Power-Sum Symmetric Functions}
\label{subsec:power-sum-symmetric}
%%%%%%%%%%%%%%%%%%%%%%%%%%%%%%%%%%%%%%%%%%%%%%%%%%%%%%%%%%%%

Define

\[
\boxed{
p_k
=
x_1^k+x_2^k+\cdots+x_n^k.
}
\]

Products

\[
p_\lambda
=
p_{\lambda_1}\cdots p_{\lambda_r}
\]

are indexed by integer partitions

\[
\lambda.
\]

Power-sum functions are particularly natural when studying cycle types
of permutations.

%%%%%%%%%%%%%%%%%%%%%%%%%%%%%%%%%%%%%%%%%%%%%%%%%%%%%%%%%%%%
\subsection{Schur Functions}
\label{subsec:schur-functions}
%%%%%%%%%%%%%%%%%%%%%%%%%%%%%%%%%%%%%%%%%%%%%%%%%%%%%%%%%%%%

Schur functions are symmetric functions indexed by partitions:

\[
\boxed{
s_\lambda.
}
\]

They can be defined combinatorially using semistandard Young tableaux.

They are central because they connect:

\[
\boxed{
\text{Young tableaux}
+
\text{symmetric functions}
+
\text{representations of }S_n.
}
\]

Schur functions form one of the most important bases of the ring of
symmetric functions.

%%%%%%%%%%%%%%%%%%%%%%%%%%%%%%%%%%%%%%%%%%%%%%%%%%%%%%%%%%%%
\subsection{Representation Theory Connection}
\label{subsec:representation-theory-combinatorics}
%%%%%%%%%%%%%%%%%%%%%%%%%%%%%%%%%%%%%%%%%%%%%%%%%%%%%%%%%%%%

The irreducible complex representations of the symmetric group

\[
S_n
\]

are indexed by partitions

\[
\lambda\vdash n.
\]

The dimension of the irreducible representation corresponding to
\(\lambda\) is

\[
\boxed{
f^\lambda,
}
\]

the number of standard Young tableaux of shape \(\lambda\).

Thus the hook-length formula also computes representation dimensions.

%%%%%%%%%%%%%%%%%%%%%%%%%%%%%%%%%%%%%%%%%%%%%%%%%%%%%%%%%%%%
\subsection{Robinson--Schensted Correspondence}
\label{subsec:robinson-schensted}
%%%%%%%%%%%%%%%%%%%%%%%%%%%%%%%%%%%%%%%%%%%%%%%%%%%%%%%%%%%%

The Robinson--Schensted correspondence associates a permutation

\[
\pi\in S_n
\]

with a pair

\[
(P,Q)
\]

of standard Young tableaux of the same shape.

Symbolically,

\[
\boxed{
\pi
\longleftrightarrow
(P,Q).
}
\]

This correspondence is bijective.

Consequently,

\[
\boxed{
n!
=
\sum_{\lambda\vdash n}
(f^\lambda)^2.
}
\]

This identity also reflects the decomposition of the regular
representation of \(S_n\).

%%%%%%%%%%%%%%%%%%%%%%%%%%%%%%%%%%%%%%%%%%%%%%%%%%%%%%%%%%%%
\subsection{Longest Increasing Subsequences}
\label{subsec:longest-increasing-subsequences}
%%%%%%%%%%%%%%%%%%%%%%%%%%%%%%%%%%%%%%%%%%%%%%%%%%%%%%%%%%%%

The Robinson--Schensted correspondence contains information about
subsequences of permutations.

The length of the first row of the associated Young diagram equals the
length of the longest increasing subsequence.

Similarly, the first column records the longest decreasing subsequence.

Thus tableau shape converts a sequence-ordering problem into a partition
problem.

%%%%%%%%%%%%%%%%%%%%%%%%%%%%%%%%%%%%%%%%%%%%%%%%%%%%%%%%%%%%
\subsection{Incidence Algebra}
\label{subsec:incidence-algebra}
%%%%%%%%%%%%%%%%%%%%%%%%%%%%%%%%%%%%%%%%%%%%%%%%%%%%%%%%%%%%

Let \(P\) be a locally finite poset.

An incidence algebra consists of functions

\[
f(x,y)
\]

defined for comparable pairs

\[
x\leq y.
\]

The product is convolution:

\[
\boxed{
(f*g)(x,z)
=
\sum_{x\leq y\leq z}
f(x,y)g(y,z).
}
\]

Incidence algebras provide an algebraic framework for inversion on
partially ordered sets.

%%%%%%%%%%%%%%%%%%%%%%%%%%%%%%%%%%%%%%%%%%%%%%%%%%%%%%%%%%%%
\subsection{Zeta Function of a Poset}
\label{subsec:zeta-function-poset}
%%%%%%%%%%%%%%%%%%%%%%%%%%%%%%%%%%%%%%%%%%%%%%%%%%%%%%%%%%%%

The zeta function of a poset is

\[
\boxed{
\zeta(x,y)
=
\begin{cases}
1,&x\leq y,\\
0,&\text{otherwise}.
\end{cases}
}
\]

It acts as a cumulative summation operator over intervals of the poset.

The inverse of \(\zeta\) under incidence-algebra convolution is the
Möbius function.

%%%%%%%%%%%%%%%%%%%%%%%%%%%%%%%%%%%%%%%%%%%%%%%%%%%%%%%%%%%%
\subsection{Möbius Function of a Poset}
\label{subsec:mobius-function-poset}
%%%%%%%%%%%%%%%%%%%%%%%%%%%%%%%%%%%%%%%%%%%%%%%%%%%%%%%%%%%%

The Möbius function

\[
\mu(x,y)
\]

is defined recursively by

\[
\boxed{
\mu(x,x)=1
}
\]

and, for \(x<y\),

\[
\boxed{
\sum_{x\leq z\leq y}
\mu(x,z)=0.
}
\]

The Greek letter

\[
\mu
\]

is called ``mu.''

The Möbius function generalizes the alternating signs found in
inclusion--exclusion.

%%%%%%%%%%%%%%%%%%%%%%%%%%%%%%%%%%%%%%%%%%%%%%%%%%%%%%%%%%%%
\subsection{Möbius Inversion}
\label{subsec:mobius-inversion-posets}
%%%%%%%%%%%%%%%%%%%%%%%%%%%%%%%%%%%%%%%%%%%%%%%%%%%%%%%%%%%%

\begin{theorem}[Möbius Inversion]
Suppose

\[
g(x)
=
\sum_{y\leq x}f(y).
\]

Then

\[
\boxed{
f(x)
=
\sum_{y\leq x}
\mu(y,x)g(y).
}
\]
\end{theorem}

Thus cumulative summation over a poset can be inverted algebraically.

Classical binomial inversion and number-theoretic Möbius inversion are
special cases of this broader principle.

%%%%%%%%%%%%%%%%%%%%%%%%%%%%%%%%%%%%%%%%%%%%%%%%%%%%%%%%%%%%
\subsection{Boolean Lattice Möbius Function}
\label{subsec:boolean-lattice-mobius}
%%%%%%%%%%%%%%%%%%%%%%%%%%%%%%%%%%%%%%%%%%%%%%%%%%%%%%%%%%%%

In the Boolean lattice

\[
\mathcal{P}([n]),
\]

if

\[
A\subseteq B,
\]

then

\[
\boxed{
\mu(A,B)
=
(-1)^{|B\setminus A|}.
}
\]

Thus inclusion--exclusion is exactly Möbius inversion on the Boolean
lattice.

This provides an algebraic explanation for alternating signs in
inclusion--exclusion formulas.

%%%%%%%%%%%%%%%%%%%%%%%%%%%%%%%%%%%%%%%%%%%%%%%%%%%%%%%%%%%%
\subsection{Chromatic Polynomial Preview}
\label{subsec:chromatic-polynomial-preview}
%%%%%%%%%%%%%%%%%%%%%%%%%%%%%%%%%%%%%%%%%%%%%%%%%%%%%%%%%%%%

For a graph \(G\), the chromatic polynomial

\[
\boxed{
P_G(q)
}
\]

counts proper vertex colorings using \(q\) available colors.

Although graph coloring will be developed later, the chromatic
polynomial is an important algebraic-combinatorial invariant.

For example, for the complete graph

\[
K_n,
\]

\[
\boxed{
P_{K_n}(q)
=
q(q-1)(q-2)\cdots(q-n+1).
}
\]

%%%%%%%%%%%%%%%%%%%%%%%%%%%%%%%%%%%%%%%%%%%%%%%%%%%%%%%%%%%%
\subsection{Deletion--Contraction}
\label{subsec:deletion-contraction-preview}
%%%%%%%%%%%%%%%%%%%%%%%%%%%%%%%%%%%%%%%%%%%%%%%%%%%%%%%%%%%%

For an edge \(e\) of a graph \(G\), the chromatic polynomial satisfies

\[
\boxed{
P_G(q)
=
P_{G-e}(q)
-
P_{G/e}(q).
}
\]

Here:

\[
G-e
\]

means delete \(e\), while

\[
G/e
\]

means contract \(e\).

Deletion--contraction is a recurring algebraic-combinatorial recurrence.

It appears again in matroid theory and the Tutte polynomial.

%%%%%%%%%%%%%%%%%%%%%%%%%%%%%%%%%%%%%%%%%%%%%%%%%%%%%%%%%%%%
\subsection{Tutte Polynomial Preview}
\label{subsec:tutte-polynomial-preview}
%%%%%%%%%%%%%%%%%%%%%%%%%%%%%%%%%%%%%%%%%%%%%%%%%%%%%%%%%%%%

The Tutte polynomial

\[
\boxed{
T_G(x,y)
}
\]

is a two-variable graph invariant satisfying deletion--contraction
relations.

Many graph quantities can be recovered from special evaluations of
\(T_G\).

Examples include information about:

\begin{itemize}
    \item spanning trees;
    \item colorings;
    \item flows;
    \item connectivity.
\end{itemize}

The Tutte polynomial becomes especially natural in matroid theory.

%%%%%%%%%%%%%%%%%%%%%%%%%%%%%%%%%%%%%%%%%%%%%%%%%%%%%%%%%%%%
\subsection{Adjacency Algebra}
\label{subsec:adjacency-algebra}
%%%%%%%%%%%%%%%%%%%%%%%%%%%%%%%%%%%%%%%%%%%%%%%%%%%%%%%%%%%%

Matrices associated with combinatorial objects may generate algebras.

For a graph with adjacency matrix \(A\), the matrices

\[
I,A,A^2,A^3,\ldots
\]

encode walk structure.

Polynomial expressions in \(A\),

\[
p(A),
\]

form an algebra under matrix addition and multiplication.

Spectral properties of \(A\) often reveal combinatorial information.

%%%%%%%%%%%%%%%%%%%%%%%%%%%%%%%%%%%%%%%%%%%%%%%%%%%%%%%%%%%%
\subsection{Eigenvalues as Combinatorial Invariants}
\label{subsec:eigenvalues-combinatorial-invariants}
%%%%%%%%%%%%%%%%%%%%%%%%%%%%%%%%%%%%%%%%%%%%%%%%%%%%%%%%%%%%

If two graphs are isomorphic, their adjacency matrices differ by a
permutation similarity transformation.

Therefore they have the same eigenvalues.

Thus the adjacency spectrum is an isomorphism invariant.

However, the converse is false:

\[
\boxed{
\text{same spectrum}
\not\Longrightarrow
\text{isomorphic graphs}.
}
\]

Graphs with the same spectrum are called cospectral.

%%%%%%%%%%%%%%%%%%%%%%%%%%%%%%%%%%%%%%%%%%%%%%%%%%%%%%%%%%%%
\subsection{Association Schemes}
\label{subsec:association-schemes}
%%%%%%%%%%%%%%%%%%%%%%%%%%%%%%%%%%%%%%%%%%%%%%%%%%%%%%%%%%%%

An association scheme partitions pairs of elements into several regular
relation classes.

It generalizes highly symmetric graph structures.

Association schemes connect:

\begin{itemize}
    \item algebraic combinatorics;
    \item coding theory;
    \item design theory;
    \item distance-regular graphs;
    \item representation theory.
\end{itemize}

Their adjacency matrices span a finite-dimensional algebra called the
Bose--Mesner algebra.

%%%%%%%%%%%%%%%%%%%%%%%%%%%%%%%%%%%%%%%%%%%%%%%%%%%%%%%%%%%%
\subsection{Bose--Mesner Algebra}
\label{subsec:bose-mesner-algebra}
%%%%%%%%%%%%%%%%%%%%%%%%%%%%%%%%%%%%%%%%%%%%%%%%%%%%%%%%%%%%

Suppose an association scheme has adjacency matrices

\[
A_0,A_1,\ldots,A_d.
\]

Their linear span

\[
\boxed{
\mathcal{A}
=
\operatorname{span}
\{A_0,\ldots,A_d\}
}
\]

is the Bose--Mesner algebra.

This algebra is closed under matrix multiplication and contains
important spectral information about the scheme.

It provides a bridge between finite combinatorial relations and linear
algebra.

%%%%%%%%%%%%%%%%%%%%%%%%%%%%%%%%%%%%%%%%%%%%%%%%%%%%%%%%%%%%
\subsection{Distance-Regular Graphs}
\label{subsec:distance-regular-graphs-preview}
%%%%%%%%%%%%%%%%%%%%%%%%%%%%%%%%%%%%%%%%%%%%%%%%%%%%%%%%%%%%

A distance-regular graph is a highly symmetric graph in which the number
of vertices at specified relative distances depends only on the
distance between two chosen vertices.

These graphs are closely related to association schemes.

They appear in:

\begin{itemize}
    \item coding theory;
    \item finite geometry;
    \item spectral graph theory;
    \item design theory.
\end{itemize}

%%%%%%%%%%%%%%%%%%%%%%%%%%%%%%%%%%%%%%%%%%%%%%%%%%%%%%%%%%%%
\subsection{Algebraic Enumeration Strategy}
\label{subsec:algebraic-enumeration-strategy}
%%%%%%%%%%%%%%%%%%%%%%%%%%%%%%%%%%%%%%%%%%%%%%%%%%%%%%%%%%%%

A common workflow in algebraic combinatorics is:

\begin{enumerate}
    \item identify the combinatorial objects;
    \item identify symmetries or algebraic structure;
    \item encode the structure using a group, polynomial, matrix, or
          symmetric function;
    \item perform algebraic operations;
    \item translate the algebraic result back into combinatorial
          information.
\end{enumerate}

This is the central methodological pattern of the subject.

%%%%%%%%%%%%%%%%%%%%%%%%%%%%%%%%%%%%%%%%%%%%%%%%%%%%%%%%%%%%
\subsection{Common Errors}
\label{subsec:algebraic-combinatorics-common-errors}
%%%%%%%%%%%%%%%%%%%%%%%%%%%%%%%%%%%%%%%%%%%%%%%%%%%%%%%%%%%%

\subsubsection{Dividing by the Group Size Without Checking Stabilizers}

The number of unlabeled objects is not generally

\[
\frac{|X|}{|G|}.
\]

Different objects may have different stabilizer sizes.

Burnside's lemma correctly handles this variation.

%%%%%%%%%%%%%%%%%%%%%%%%%%%%%%%%%%%%%%%%%%%%%%%%%%%%%%%%%%%%
\subsubsection{Confusing Orbits and Stabilizers}

The orbit describes where an object can move.

The stabilizer describes which symmetries leave it fixed.

They are connected by

\[
|G|
=
|\operatorname{Orb}(x)|
|\operatorname{Stab}(x)|.
\]

%%%%%%%%%%%%%%%%%%%%%%%%%%%%%%%%%%%%%%%%%%%%%%%%%%%%%%%%%%%%
\subsubsection{Counting Fixed Points Incorrectly}

When a permutation of positions has several cycles, a fixed coloring
must be constant on each cycle.

With \(q\) colors, the count is

\[
q^{c(g)},
\]

not necessarily \(q^n\).

%%%%%%%%%%%%%%%%%%%%%%%%%%%%%%%%%%%%%%%%%%%%%%%%%%%%%%%%%%%%
\subsubsection{Confusing Burnside and P\'olya Enumeration}

Burnside's lemma counts orbits.

P\'olya enumeration refines the count by tracking color multiplicities
using the cycle index.

%%%%%%%%%%%%%%%%%%%%%%%%%%%%%%%%%%%%%%%%%%%%%%%%%%%%%%%%%%%%
\subsubsection{Confusing Ordinary and Gaussian Binomial Coefficients}

The ordinary coefficient

\[
\binom{n}{k}
\]

counts \(k\)-element subsets.

The Gaussian coefficient

\[
\begin{bmatrix}
n\\
k
\end{bmatrix}_q
\]

counts \(k\)-dimensional subspaces of \(\mathbb{F}_q^n\) when \(q\) is
a prime power.

%%%%%%%%%%%%%%%%%%%%%%%%%%%%%%%%%%%%%%%%%%%%%%%%%%%%%%%%%%%%
\subsubsection{Forgetting What \(q\) Represents}

In some contexts, \(q\) is a formal variable.

In finite-field counting, \(q\) is a prime power.

The interpretation depends on context.

%%%%%%%%%%%%%%%%%%%%%%%%%%%%%%%%%%%%%%%%%%%%%%%%%%%%%%%%%%%%
\subsubsection{Confusing Integer Partitions with Set Partitions}

Young diagrams are indexed by integer partitions, not set partitions.

%%%%%%%%%%%%%%%%%%%%%%%%%%%%%%%%%%%%%%%%%%%%%%%%%%%%%%%%%%%%
\subsubsection{Miscomputing Hook Lengths}

The hook of a cell includes:

\[
\boxed{
\text{the cell itself}
+
\text{cells to the right}
+
\text{cells below}.
}
\]

Diagonal cells are not included unless they lie directly right or below.

%%%%%%%%%%%%%%%%%%%%%%%%%%%%%%%%%%%%%%%%%%%%%%%%%%%%%%%%%%%%
\subsubsection{Assuming Spectral Equality Implies Isomorphism}

Isomorphic graphs are cospectral.

Cospectral graphs need not be isomorphic.

%%%%%%%%%%%%%%%%%%%%%%%%%%%%%%%%%%%%%%%%%%%%%%%%%%%%%%%%%%%%
\subsection{Applications}
\label{subsec:algebraic-combinatorics-applications}
%%%%%%%%%%%%%%%%%%%%%%%%%%%%%%%%%%%%%%%%%%%%%%%%%%%%%%%%%%%%

\begin{applicationbox}

\textbf{Symmetry Enumeration.}
Burnside's lemma and P\'olya enumeration count necklaces, bracelets,
tilings, colorings, and molecular substitution patterns.

\medskip

\textbf{Representation Theory.}
Young tableaux and Schur functions describe irreducible representations
of symmetric groups.

\medskip

\textbf{Coding Theory.}
Gaussian binomial coefficients, association schemes, and
distance-regular graphs arise naturally in finite coding spaces.

\medskip

\textbf{Graph Theory.}
Chromatic polynomials, Tutte polynomials, adjacency spectra, and
association schemes encode graph structure algebraically.

\medskip

\textbf{Finite Geometry.}
Subspace counting over finite fields produces Gaussian binomial
coefficients.

\medskip

\textbf{Chemistry.}
Molecular symmetry groups permit systematic enumeration of distinct
substitution patterns.

\medskip

\textbf{Statistical Physics.}
Symmetry reduction prevents equivalent configurations from being
counted repeatedly.

\medskip

\textbf{Computer Science.}
Canonical forms, permutation groups, symmetry reduction, and graph
isomorphism algorithms rely on algebraic-combinatorial structure.

\end{applicationbox}

%%%%%%%%%%%%%%%%%%%%%%%%%%%%%%%%%%%%%%%%%%%%%%%%%%%%%%%%%%%%
\subsection{Exercises}
\label{subsec:algebraic-combinatorics-exercises}
%%%%%%%%%%%%%%%%%%%%%%%%%%%%%%%%%%%%%%%%%%%%%%%%%%%%%%%%%%%%

\begin{exercise}[1]
Define the orbit of an element under a group action.
\end{exercise}

\begin{exercise}[1]
Define the stabilizer of an element under a group action.
\end{exercise}

\begin{exercise}[1]
State the orbit--stabilizer theorem.
\end{exercise}

\begin{exercise}[1]
State Burnside's lemma.
\end{exercise}

\begin{exercise}[1]
A symmetry group has order \(12\), and an object has stabilizer of size
\(3\). Find the size of its orbit.
\end{exercise}

\begin{exercise}[1]
How many binary colorings of four labeled positions exist before
symmetry is considered?
\end{exercise}

\begin{exercise}[1]
Compute

\[
[4]_q.
\]
\end{exercise}

\begin{exercise}[1]
Compute

\[
[3]_q!.
\]
\end{exercise}

\begin{exercise}[1]
Find

\[
\begin{bmatrix}
3\\
1
\end{bmatrix}_q.
\]
\end{exercise}

\begin{exercise}[1]
Draw the Young diagram of

\[
\lambda=(5,3,1).
\]
\end{exercise}

\begin{exercise}[2]
Use Burnside's lemma to count binary necklaces of length \(3\).
\end{exercise}

\begin{exercise}[2]
Use Burnside's lemma to count binary necklaces of length \(5\).
\end{exercise}

\begin{exercise}[2]
Count colorings of the vertices of a square using two colors up to
rotation.
\end{exercise}

\begin{exercise}[2]
Repeat the previous problem when reflections are also considered
equivalent.
\end{exercise}

\begin{exercise}[2]
Find the cycle index of \(C_3\) acting on three positions.
\end{exercise}

\begin{exercise}[2]
Use that cycle index to count colorings with \(q\) colors.
\end{exercise}

\begin{exercise}[2]
Show that a permutation with \(c\) cycles fixes exactly \(q^c\)
colorings of its positions using \(q\) colors.
\end{exercise}

\begin{exercise}[2]
Compute

\[
\begin{bmatrix}
4\\
2
\end{bmatrix}_q.
\]
\end{exercise}

\begin{exercise}[2]
Evaluate

\[
\begin{bmatrix}
4\\
2
\end{bmatrix}_q
\]

at \(q=1\).
\end{exercise}

\begin{exercise}[2]
How many one-dimensional subspaces does

\[
\mathbb{F}_2^3
\]

contain?
\end{exercise}

\begin{exercise}[2]
Use the hook-length formula to find the number of standard Young
tableaux of shape

\[
(3,1).
\]
\end{exercise}

\begin{exercise}[2]
Write the first three elementary symmetric polynomials in

\[
x_1,x_2,x_3.
\]
\end{exercise}

\begin{exercise}[2]
Show that

\[
\prod_{i=1}^{n}(1+x_it)
=
\sum_{k=0}^{n}e_kt^k.
\]
\end{exercise}

\begin{exercise}[3]
Prove Burnside's lemma by double counting pairs

\[
(g,x)
\]

with \(g\cdot x=x\).
\end{exercise}

\begin{exercise}[3]
Derive the necklace formula

\[
N_q(n)
=
\frac1n
\sum_{d\mid n}
\varphi(d)q^{n/d}.
\]
\end{exercise}

\begin{exercise}[3]
Derive the cycle index of \(C_4\).
\end{exercise}

\begin{exercise}[3]
Use P\'olya enumeration to count rotationally distinct colorings of a
square with exactly two red and two blue vertices.
\end{exercise}

\begin{exercise}[3]
Derive the \(q\)-Pascal recurrence for Gaussian binomial coefficients.
\end{exercise}

\begin{exercise}[3]
Prove that

\[
\begin{bmatrix}
n\\
k
\end{bmatrix}_q
\]

counts \(k\)-dimensional subspaces of

\[
\mathbb{F}_q^n.
\]
\end{exercise}

\begin{exercise}[3]
Prove

\[
[n]_q!
=
\sum_{\pi\in S_n}
q^{\operatorname{inv}(\pi)}.
\]
\end{exercise}

\begin{exercise}[3]
Compute the hook lengths of every cell in the Young diagram of

\[
(3,2).
\]
\end{exercise}

\begin{exercise}[3]
Use the hook-length formula to find

\[
f^{(3,2)}.
\]
\end{exercise}

\begin{exercise}[3]
Verify for \(n=3\) that

\[
3!
=
\sum_{\lambda\vdash3}
(f^\lambda)^2.
\]
\end{exercise}

\begin{exercise}[3]
Show that inclusion--exclusion is Möbius inversion on the Boolean
lattice.
\end{exercise}

\begin{exercise}[3]
Find the Möbius function of a finite chain.
\end{exercise}

\begin{exercise}[3]
Compute the chromatic polynomial of a tree with \(n\) vertices.
\end{exercise}

\begin{exercise}[3]
Verify deletion--contraction for a triangle.
\end{exercise}

\begin{exercise}[4]
Study P\'olya's enumeration theorem in full generality.
\end{exercise}

\begin{exercise}[4]
Derive the cycle index polynomial of the dihedral group \(D_n\).
\end{exercise}

\begin{exercise}[4]
Use P\'olya enumeration to count bracelets with a prescribed color
inventory.
\end{exercise}

\begin{exercise}[4]
Investigate the relationship between Gaussian binomial coefficients and
partitions fitting inside a rectangle.
\end{exercise}

\begin{exercise}[4]
Study Schur functions and their definition through semistandard Young
tableaux.
\end{exercise}

\begin{exercise}[4]
Investigate the Robinson--Schensted correspondence.
\end{exercise}

\begin{exercise}[4]
Study the relationship between longest increasing subsequences and the
shape produced by Robinson--Schensted.
\end{exercise}

\begin{exercise}[4]
Investigate incidence algebras and Möbius inversion on general locally
finite posets.
\end{exercise}

\begin{exercise}[4]
Study the Tutte polynomial and its deletion--contraction recurrence.
\end{exercise}

\begin{exercise}[4]
Investigate association schemes and the Bose--Mesner algebra.
\end{exercise}

%%%%%%%%%%%%%%%%%%%%%%%%%%%%%%%%%%%%%%%%%%%%%%%%%%%%%%%%%%%%
\subsection{Conceptual Summary}
\label{subsec:algebraic-combinatorics-summary}
%%%%%%%%%%%%%%%%%%%%%%%%%%%%%%%%%%%%%%%%%%%%%%%%%%%%%%%%%%%%

Algebraic combinatorics uses algebraic structure to study discrete
objects.

A group action organizes objects into orbits.

For \(x\in X\),

\[
\boxed{
|G|
=
|\operatorname{Orb}(x)|
|\operatorname{Stab}(x)|.
}
\]

Burnside's lemma counts symmetry classes:

\[
\boxed{
|X/G|
=
\frac1{|G|}
\sum_{g\in G}
|\operatorname{Fix}(g)|.
}
\]

For colorings, a symmetry with \(c(g)\) cycles fixes

\[
\boxed{
q^{c(g)}
}
\]

colorings.

Cycle index polynomials organize cycle structure, while P\'olya
enumeration refines orbit counts by color content.

The \(q\)-integer is

\[
\boxed{
[n]_q
=
1+q+\cdots+q^{n-1}.
}
\]

Gaussian binomial coefficients

\[
\boxed{
\begin{bmatrix}
n\\
k
\end{bmatrix}_q
}
\]

count \(k\)-dimensional subspaces of \(\mathbb{F}_q^n\).

Young diagrams encode integer partitions.

The hook-length formula gives

\[
\boxed{
f^\lambda
=
\frac{n!}{
\prod_{u\in\lambda}h(u)
}.
}
\]

Symmetric functions and Schur functions connect tableaux with
representation theory.

Möbius inversion generalizes inclusion--exclusion to partially ordered
sets.

Polynomial and matrix invariants such as chromatic polynomials,
Tutte polynomials, and adjacency spectra translate combinatorial
structure into algebra.

The overall theme is

\[
\boxed{
\text{symmetry}
+
\text{algebra}
+
\text{enumeration}
+
\text{structure}.
}
\]

%%%%%%%%%%%%%%%%%%%%%%%%%%%%%%%%%%%%%%%%%%%%%%%%%%%%%%%%%%%%
\subsection{Section Connections}
\label{subsec:algebraic-combinatorics-connections}
%%%%%%%%%%%%%%%%%%%%%%%%%%%%%%%%%%%%%%%%%%%%%%%%%%%%%%%%%%%%

\chapterconnection{
Algebraic Combinatorics builds on the symmetry, permutation, partition,
and generating-function ideas developed throughout the earlier sections
of Chapter 6. It connects especially strongly with Chapter~\ref{chap:algebra}
through group actions, symmetric groups, representation theory, finite
fields, and polynomial algebra. The next section, Probabilistic
Combinatorics, replaces algebraic averaging over symmetry groups with
probabilistic averaging over random structures. Graph Theory later
develops chromatic, spectral, and structural graph concepts in greater
detail. Design Theory uses group actions, incidence algebras, and
association schemes, while Matroid Theory generalizes deletion--
contraction and Tutte-polynomial methods to abstract independence
structures.
}

%%%%%%%%%%%%%%%%%%%%%%%%%%%%%%%%%%%%%%%%%%%%%%%%%%%%%%%%%%%%
% SECTION SOURCE TRACKING
%%%%%%%%%%%%%%%%%%%%%%%%%%%%%%%%%%%%%%%%%%%%%%%%%%%%%%%%%%%%

%------------------------------------------------------------
% 6.7 Algebraic Combinatorics
%------------------------------------------------------------
%
% Source basis:
%   - Standard algebraic combinatorics.
%   - Standard finite group actions.
%   - Standard symmetry enumeration.
%   - Standard symmetric-function and tableau theory.
%   - Standard incidence-algebra methods.
%
% Used for:
%   - group actions;
%   - orbits and stabilizers;
%   - orbit--stabilizer theorem;
%   - Burnside's lemma;
%   - fixed-point enumeration;
%   - cycle structure;
%   - necklace and bracelet counting;
%   - cycle index polynomials;
%   - Polya enumeration;
%   - color-content enumeration;
%   - permutation statistics;
%   - q-integers;
%   - q-factorials;
%   - Gaussian binomial coefficients;
%   - finite-field subspace counting;
%   - Young diagrams;
%   - Young tableaux;
%   - hook lengths;
%   - hook-length formula;
%   - Young's lattice;
%   - symmetric polynomials;
%   - elementary symmetric functions;
%   - complete homogeneous symmetric functions;
%   - power-sum symmetric functions;
%   - Schur functions;
%   - representation-theoretic connections;
%   - Robinson--Schensted correspondence;
%   - longest increasing subsequences;
%   - incidence algebras;
%   - zeta and Möbius functions of posets;
%   - Möbius inversion;
%   - Boolean-lattice inclusion--exclusion;
%   - chromatic polynomial preview;
%   - deletion--contraction;
%   - Tutte polynomial preview;
%   - adjacency algebra;
%   - spectral invariants;
%   - association schemes;
%   - Bose--Mesner algebra;
%   - distance-regular graphs;
%   - applications and exercises.
%
% Notes:
%   - Group theory and representation theory are developed more
%     systematically in Chapter 1.
%   - Graph polynomial and spectral ideas are revisited in Graph Theory.
%   - Association schemes connect strongly with Design Theory.
%   - Deletion--contraction and the Tutte polynomial reappear naturally
%     in Matroid Theory.
%
%%%%%%%%%%%%%%%%%%%%%%%%%%%%%%%%%%%%%%%%%%%%%%%%%%%%%%%%%%%%